\addtocontents{toc}{\protect\Needspace{3\baselineskip}}
\originalchapter{宇宙热史与轻元素}\label{ch:thermal}
\originalsection{相对论粒子与熵}
温度高于静止质量时, 粒子速度接近 $c$. 对自由玻色子与费米子, 占据数分别为 $(\ee^{E/T}-1)^{-1}$ 和 $(\ee^{E/T}+1)^{-1}$; 本节恢复 $\kb,\hbar,c$ 时写明单位. 动量积分给
\begin{equation}
u=\frac{\pi^2}{30}g_*\frac{(\kb T)^4}{(\hbar c)^3},\qquad
s=\frac{2\pi^2}{45}g_{*s}\kb\frac{(\kb T)^3}{(\hbar c)^3}.
\end{equation}
费米子每自由度对 $g_*$ 的权重为 $7/8$, 若不同粒子温度不同还乘 $(T_i/T)^4$; 熵权重则乘三次方. 辐射压强 $p=u/3$ 由各向同性动量流得到, $s=(u+p)/T$ 使用零化学势. 绝热膨胀时 $sa^3$ 守恒, 所以 $aT\propto g_{*s}^{-1/3}$. 粒子湮灭时 $g_{*s}$ 下降, 仍耦合的粒子升温得比简单 $a^{-1}$ 慢.

\originalsection{反应率与冻结}
是否保持平衡取决于每粒子反应率 $\Gamma=n\langle\sigma v\rangle$ 和膨胀率 $H$, 而不只是温度高低. $\Gamma\gg H$ 时体系能跟随平衡分布; $\Gamma\lesssim H$ 时反应来不及调整. 弱相互作用中微子粗略 $\Gamma\sim G_F^2T^5$, 辐射期 $H\sim T^2/\mpl$, 故宇宙冷却时前者下降更快, 在 MeV 尺度退耦. 此估计不包含精确输运和非瞬时退耦修正.

中微子先退耦后, 电子正电子湮灭只加热光子. 湮灭前与光子耦合的熵自由度为 $2+(7/8)4=11/2$, 后为 2, 因而 $T_\gamma/T_\nu=(11/4)^{1/3}$. 中微子温度继续按 $a^{-1}$ 下降. 暗物质若以热遗迹产生, 数密度满足 $\dot n+3Hn=-\langle\sigma v\rangle(n^2-n_{\rm eq}^2)$. 用 $Y=n/s$ 消去纯稀释后, 剩余丰度取决于冻结期间的反应积分, 不能只给 $n=n_{\rm eq}$ 并认为今天仍平衡.

\originalsection{核合成与复合}
弱平衡阶段的中子质子比近似 $n_n/n_p\simeq\exp[-(m_n-m_p)c^2/\kb T]$. 弱反应冻结后中子继续衰变, 直到氘能存活并迅速生成氦. 由于每重子有许多光子, 即使 $\kb T$ 已低于氘结合能, 高频尾仍能毁氘, 形成氘瓶颈. 当核反应展开, 大部分存活中子进入 ${}^4\mathrm{He}$, 故氦质量分数近似 $Y_p\simeq2r/(1+r)$, $r$ 是该时刻中子质子比. 精确 D、He、Li 丰度需要反应网络和速率数据; 此处不把估计标成精密 BBN 计算.

氢复合也发生在远低于 $13.6\,\mathrm{eV}$ 的温度, 原因同样与光子重子比有关. Saha 关系描述早期平衡, 但后期非平衡复合需要速率方程. 光子的 Thomson 散射率 $n_e\sigma_Tc$ 随自由电子减少而下降. 定义向今日的光学深度 $\tau$, 最后散射概率密度与 $\ee^{-\tau}|\dot\tau|$ 成正比. 复合、光子退耦与后来再电离是不同过程, 不能视为同一瞬间.
